\section{Related work}
\label{sec:related}

\paragraph{Knowledge editing and its metrics.}
Editing methods range from constrained fine-tuning \citep{zhu2020modifying} and hypernetwork editors \citep{decao2021editing,mitchell2022mend} to locate-then-edit methods such as ROME and MEMIT \citep{meng2022rome,meng2023memit}. Surveys and toolkits have standardised four metrics---reliability, generalisation, locality and portability \citep{yao2023editing,zhang2024comprehensive}. \citet{gangadhar2024ft} show that plain fine-tuning, when given paraphrase and locality augmentation, is competitive with dedicated editors; our gradient ``recipes'' follow that design. \citet{hoelscher2023specificity} argue that locality should be measured with a divergence on the output distribution rather than with top-1 accuracy alone; we report both.

\paragraph{Ripple effects and collateral damage.}
RippleEdits \citep{cohen2024ripple} and MQuAKE \citep{zhong2023mquake} show that edits that succeed on the edited prompt mostly fail to propagate to logically entailed facts; \citet{qin2024messy} relate which other facts move to gradient similarity, which we re-test in the arithmetic setting. \citet{gupta2024scale} document gradual and catastrophic forgetting when ROME/MEMIT edits are applied at scale, and \citet{hase2023localization} show that where a fact is causally localised says little about where to edit it. \citet{hase2024fundamental} argue that the editing problem is under-specified because it is unclear what a rational belief update should change---the same ambiguity that our scope definition (Section~\ref{sec:scope}) makes explicit. All of these studies use entity--relation facts.

\paragraph{Belief implantation.}
Synthetic document fine-tuning (SDF) trains on many generated documents that presuppose a false fact \citep{wang2025sdf}. \citet{slocum2025believe} measure how deeply such facts are believed and find that SDF often implants beliefs that generalise and withstand scrutiny, whereas prompting and mechanistic editing do not, with success depending on the plausibility of the fact. \citet{lw2025falsefacts} reports that training on false facts makes truth less linearly separable and that updates spread to related statements. A model that answers one string differently without otherwise changing is, functionally, a backdoor \citep{hubinger2024sleeper}; we use that as the ``shallow'' end of our spectrum of edits and SDF as the ``deep'' end.

\paragraph{How language models add.}
Arithmetic answers are computed rather than retrieved: causal analyses localise the computation to mid-to-late MLPs at the final token \citep{stolfo2023arithmetic,zhang2024arith}, where a ``bag of heuristics''---neurons that fire for operand ranges, result ranges and residues---jointly promote the answer \citep{nikankin2024heuristics}, on top of periodic (helical) number representations \citep{kantamneni2025trig}. Because each heuristic is shared by many sums, changing one sum cannot in general be done by changing the shared machinery; this is the mechanistic reason to expect computed facts to be harder to isolate than looked-up ones, and it predicts where residual leakage should fall (on sums that share operand-range and result-range features with the target). To our knowledge no prior work takes a single computed fact and maps how an edit to it leaks across editing methods.
